\chapter{面向微调鲁棒性的双子空间约束安全知识编辑方法}\label{chap:method}

本章研究安全知识编辑在后续微调作用下的保持问题。现有参数级编辑通常在有限编辑样本上求得一次参数更新；当模型随后进行领域适配、指令跟随或其他任务微调时，新的梯度更新可能改变编辑所依赖的参数方向，使安全响应减弱。与此同时，直接增大安全编辑幅度可能损害模型对无害请求的正常响应。为此，本章提出一种双子空间约束的鲁棒安全编辑方案：利用安全编辑样本和无害任务样本的梯度分别构造安全子空间与效用保护子空间，将编辑更新限制在安全子空间内，并抑制其在效用保护方向上的分量；进一步通过代表性后续微调扰动构造鲁棒目标，使编辑阶段能够显式考虑微调后的安全行为。

本章将先给出问题定义、威胁模型和评价指标，再说明子空间构建、联合约束及微调扰动鲁棒优化，最后给出实验方案。由于本稿尚未包含实际实验记录，本章的实验结果、图表数值和基于结果的结论均以待补项标明。

\section{相关知识介绍}

\subsection{参数级知识编辑}

知识编辑旨在以有限计算代价改变模型在特定输入上的输出，而不对模型进行全量重新训练。设预训练模型为$f_{\theta_0}$，参数为$\theta_0$，编辑请求集合为$\mathcal{D}_e=\{(x_i,y_i^*)\}_{i=1}^{N_e}$，其中$x_i$为输入，$y_i^*$为期望输出。编辑器根据$\mathcal{D}_e$产生参数增量$\delta$，得到编辑后模型$f_{\theta_0+\delta}$。一个有效的编辑通常需要同时满足目标行为改变、相近输入上的泛化，以及无关请求上的行为保持等要求。ROME、MEMIT等方法直接更新模型的部分参数；AlphaEdit进一步将更新投影至特定知识表示的零空间，以减轻对保留知识的影响\cite{meng2022rome,meng2023memit,fang2024alphaedit}。相关研究也指出，编辑方法在定位方式、参数范围、编辑规模和知识保持目标上存在差异\cite{zhang2024comprehensive,wang2024survey}。

本文关注的不是将某条事实替换为另一条事实，而是对给定风险请求形成稳定的安全响应。对于有害请求$x$，期望输出$y^{\mathrm{safe}}$可以是符合策略的拒答，也可以是安全替代信息或求助建议；目标输出由数据集标注和安全策略确定。本文以交叉熵形式表示编辑目标：
\begin{equation}
  \mathcal{L}_{s}(\theta;\mathcal{D}_{s})
  =-\frac{1}{|\mathcal{D}_{s}|}
  \sum_{(x,y^{\mathrm{safe}})\in\mathcal{D}_{s}}
  \log p_{\theta}(y^{\mathrm{safe}}\mid x),
  \label{eq:safety-edit-loss}
\end{equation}
其中$\mathcal{D}_{s}$为安全编辑训练集。该损失使模型提高生成目标安全响应的概率，但损失下降本身并不能充分证明模型对未见攻击、改写请求或后续微调保持安全，因此仍需独立测试集和后续微调实验。

\subsection{安全知识编辑与持久性}

安全知识编辑是以改变模型对特定风险请求的响应行为为目标的参数编辑过程。SafeEdit等基准覆盖多类不安全行为及不同攻击形式，为安全编辑效果和模型能力影响提供了系统评测基础\cite{wang2024safeedit}。在该任务中，安全效果与正常效用是并列目标：模型应减少对有害请求的危险响应，同时继续回答合理的无害请求。若只报告拒答率，可能把无差别拒答误判为安全改善；若只报告一般任务准确率，也无法判断有害响应是否被有效抑制。

本文区分编辑即时效果与微调后的保持效果。令$\theta_e=\theta_0+\delta$为编辑后参数，对环境$e_j$定义$\widehat{\theta}_{\mathrm{ft}}^{(j)}=\mathrm{FT}_{e_j}(\theta_e)$，其中$\mathrm{FT}_{e_j}$表示从当前模型出发执行指定微调过程。即时效果在$f_{\theta_e}$上测量；持久性则在$f_{\widehat{\theta}_{\mathrm{ft}}^{(j)}}$上测量。编辑持久性不等于模型内部知识被彻底删除，也不意味着对任意后续训练都有保证。它表示在本文定义的微调环境、训练预算和安全测试集下，编辑行为能够保持到何种程度。关于一般知识编辑结果受到后续微调影响的实证研究，以及面向遗忘任务的微调鲁棒方法，为该问题提供了研究依据\cite{cheng2025finetune,wang2025ilu}。

\subsection{微调扰动与威胁模型}

本文考虑模型完成安全编辑后，继续接受下游任务微调的情形。微调的目的可以是领域适配、任务指令跟随或其他正常能力更新；本文不假设微调者有意绕过安全策略。设后续微调环境集合为$\mathcal{E}=\{e_0,e_1,\ldots,e_M\}$，其中$e_0$表示不继续微调，其余环境对应不同的微调数据、参数更新方式或训练预算。对环境$e_j$，在基础模型上执行预先规定的微调过程，得到参数$\theta_{\mathrm{ft}}^{(j)}$，并定义扰动
\begin{equation}
  \Delta_{0}^{(j)}
  =\mathrm{FT}_{e_j}(\theta_0)-\theta_0,\qquad
  \Delta_{0}^{(0)}=\mathbf{0}.
  \label{eq:ft-perturbation}
\end{equation}
从编辑模型重新执行微调所得到的真实参数变化记为
\begin{equation}
  \widehat{\Delta}^{(j)}(\delta)
  =\mathrm{FT}_{e_j}(\theta_0+\delta)-(\theta_0+\delta),
  \qquad
  \widehat{\theta}_{\mathrm{ft}}^{(j)}
  =\theta_0+\delta+\widehat{\Delta}^{(j)}(\delta).
  \label{eq:actual-ft-perturbation}
\end{equation}
在本章的鲁棒编辑目标中，$\Delta_{0}^{(j)}$作为从基础模型获得的代表性后续更新扰动模板；$\theta_0+\delta+\Delta_{0}^{(j)}$是用于编辑优化的虚拟扰动模型。由于从$\theta_0$获得的固定参数差不等同于从编辑模型$\theta_0+\delta$重新执行微调得到的真实参数差，本文将该模板视为训练阶段的近似，而不将其解释为精确微调轨迹。正式实验需说明扰动模板是由全参数微调还是参数高效微调得到，并给出微调数据、优化器、学习率、步数或轮数等配置；最终鲁棒性评估则须从编辑模型重新执行实际微调。

该威胁模型的边界包括：其一，后续微调环境由实验中明确列出的任务和训练设置构成，不覆盖任意未知训练过程；其二，研究对象是模型输出层面的安全响应保持，而非对模型内部知识进行不可逆删除；其三，模型架构、分词器和推理模板保持一致。若后续微调改变模型结构、替换大部分参数或采用与训练环境显著不同的目标，则需单独评估，不能由本章结果直接外推。

\subsection{安全性与效用评价指标}

为避免单一指标掩盖安全与效用之间的权衡，本文分别报告有害请求处理、后续微调保持、无害任务能力和局部性指标。设$\mathcal{T}_{s}$为安全测试集，$\mathcal{T}_{u}$为无害效用测试集，$A(f_\theta(x))$为安全判别器或经人工复核的输出标签。

安全攻击成功率（Attack Success Rate，ASR）定义为模型对有害请求生成违反安全策略内容的比例：
\begin{equation}
  \mathrm{ASR}(\theta)=
  \frac{1}{|\mathcal{T}_{s}|}
  \sum_{x\in\mathcal{T}_{s}}
  \mathbb{I}\!\left[A(f_\theta(x))=\mathrm{unsafe}\right].
  \label{eq:asr}
\end{equation}
安全响应率（Safe Response Rate，SRR）定义为$1-\mathrm{ASR}$。实际使用时应说明判别器、人工复核规则、生成长度和解码参数，并抽样检查自动评估的误判。

微调后安全保持率以编辑模型和微调后模型的安全响应率之比表示：
\begin{equation}
  \mathrm{Ret}^{(j)}
  =\frac{\mathrm{SRR}(\widehat{\theta}_{\mathrm{ft}}^{(j)})}
  {\mathrm{SRR}(\theta_0+\delta)+\epsilon},
  \label{eq:retention}
\end{equation}
其中$\epsilon$为避免分母为零的极小常数。保持率应与编辑前后的绝对SRR、ASR变化同时报告，以免即时编辑效果较弱时比值产生误导。

无害任务效用通过任务准确率、指令跟随评分或其他与所选基准匹配的指标衡量，并比较基础模型、编辑模型及微调后模型的变化。局部性用于度量编辑对无关请求分布的影响，可报告无害测试集上的任务性能差异及输出分布变化，例如：
\begin{equation}
  \mathrm{KL}_{u}
  =\frac{1}{|\mathcal{T}_{u}|}
  \sum_{x\in\mathcal{T}_{u}}
  D_{\mathrm{KL}}\!\left(
  p_{\theta_0}(\cdot\mid x)\,\Vert\,
  p_{\theta_0+\delta}(\cdot\mid x)
  \right).
  \label{eq:locality}
\end{equation}
表\ref{tab:ch2-metrics}汇总本章指标。具体使用的任务数据、判别器和聚合方式将在实验设置中说明。

\begin{table}[htbp]
  \centering
  \caption{第二章主要评价指标}
  \label{tab:ch2-metrics}
  \begin{tabular}{p{0.18\textwidth}p{0.26\textwidth}p{0.44\textwidth}}
    \hline
    评价维度 & 指标 & 含义 \\ \hline
    安全效果 & ASR、SRR & 测量有害请求被不安全响应或安全处理的比例 \\ 
    微调保持 & $\mathrm{Ret}^{(j)}$、SRR变化 & 测量不同后续微调环境下的安全行为保持情况 \\
    通用效用 & 任务准确率或任务评分 & 测量无害任务能力及编辑、微调前后的变化 \\
    编辑局部性 & $\mathrm{KL}_{u}$、无关任务差值 & 测量无关输入上的输出分布和任务性能变化 \\
    \hline
  \end{tabular}
\end{table}

\section{方案框架}

\subsection{微调衰减问题分析}

设安全编辑在参数$\theta_0$上得到更新$\delta$。若后续微调产生的更新方向与$\delta$存在较大相互作用，则模型在安全测试请求上的输出可能偏离编辑时所优化的目标。对安全损失在$\theta_0+\delta$处作一阶近似，有
\begin{equation}
  \mathcal{L}_{s}(\theta_e+\widehat{\Delta}^{(j)}(\delta))
  \approx \mathcal{L}_{s}(\theta_0+\delta)
  +\nabla_{\theta}\mathcal{L}_{s}(\theta_e)^{\mathsf T}\widehat{\Delta}^{(j)}(\delta).
  \label{eq:first-order-decay}
\end{equation}
式\eqref{eq:first-order-decay}说明，当后续微调扰动在安全损失的敏感方向上具有较大分量时，安全损失可能上升。该近似只用于解释参数变化与损失变化的局部联系，不构成对实际微调行为的保证；非线性、高阶效应和样本分布变化仍需通过完整模型推理进行评估。

本研究拟先进行小规模预实验：对未编辑模型执行安全编辑，并使用预先选定的下游任务微调；在编辑前、编辑后和微调后分别测量同一安全测试集上的ASR与SRR，同时记录下游任务效用。若结果显示后续微调后安全指标发生变化，再进一步按微调步数或训练预算绘制变化曲线。当前尚无预实验记录，故本节不报告衰减幅度或统计结论。

安全编辑还存在第二个问题：优化目标若只关注有害请求，可能连带改变无害请求的响应。本文把正常任务梯度汇集为效用保护方向，在参数更新时限制其影响；同时把后续微调造成的扰动纳入编辑目标。该设计分别处理编辑方向选择、无害效用保护和微调后安全保持三个环节。

% TODO: 补入预实验模型、数据、微调配置，以及编辑前/编辑后/微调后的ASR、SRR和效用结果；增加相应曲线或表格。

\subsection{基于安全梯度的安全子空间构建}

令$\Theta_e$表示允许编辑的参数集合，$\theta_0^{e}$为这些参数在初始模型中的向量化表示，维度为$d_e$。下文中的梯度和子空间均定义在$\Theta_e$对应的坐标上；为简化记号，将受限更新写作$\delta$，当其出现在完整模型表达式$f_{\theta_0+\delta}$中时，表示仅更新$\Theta_e$内的坐标，其他参数保持不变。给定安全编辑样本$\mathcal{D}_{s}^{\mathrm{train}}$，在初始模型$\theta_0$处对每个样本计算安全目标梯度：
\begin{equation}
  g_i^{s}=
  \nabla_{\theta_e}
  \ell\!\left(f_{\theta_0}(x_i),y_i^{\mathrm{safe}}\right),
  \quad (x_i,y_i^{\mathrm{safe}})\in\mathcal{D}_{s}^{\mathrm{train}}.
  \label{eq:safety-gradient}
\end{equation}
将梯度按列组成矩阵$G_s=[g_1^s,\ldots,g_{N_s}^s]\in\mathbb{R}^{d_e\times N_s}$，并对其进行截断奇异值分解：
\begin{equation}
  G_s=U_s\Sigma_sV_s^{\mathsf T},\qquad
  Q_s=U_s[:,1:r_s],\qquad P_s=Q_sQ_s^{\mathsf T}.
  \label{eq:safety-subspace}
\end{equation}
其中$Q_s$为安全子空间的正交基，$P_s$为其投影矩阵，$r_s$为子空间维度。$r_s$可依据奇异值能量累计比例确定：
\begin{equation}
  r_s=\min\left\{r:
  \frac{\sum_{k=1}^{r}\sigma_{s,k}^{2}}
  {\sum_{k}\sigma_{s,k}^{2}}\ge \rho_s\right\},
  \label{eq:safety-rank}
\end{equation}
其中$\rho_s\in(0,1]$为预先设定的能量保留阈值，并在参数敏感性实验中分析其影响。

该子空间表示训练安全样本梯度所覆盖的局部更新方向。它不表示模型所有安全行为的完整空间，也不能单独保证对分布外提示的安全性。为减少样本重复或单一提示模板对基的支配，构造$G_s$时应进行样本去重、提示改写覆盖和梯度归一化；训练、验证和测试安全样本应按提示族或语义簇划分，避免近重复样本跨集合泄漏。

\subsection{基于正常任务梯度的效用保护子空间构建}

令$\mathcal{D}_{u}^{\mathrm{train}}$为与安全编辑集分离的无害任务样本。对其任务损失计算初始模型梯度：
\begin{equation}
  g_j^{u}=
  \nabla_{\theta_e}
  \ell\!\left(f_{\theta_0}(x_j^{u}),y_j^{u}\right),
  \quad (x_j^{u},y_j^{u})\in\mathcal{D}_{u}^{\mathrm{train}}.
  \label{eq:utility-gradient}
\end{equation}
同样地，令$G_u=[g_1^u,\ldots,g_{N_u}^u]$，对$G_u$进行截断奇异值分解，并取前$r_u$个左奇异向量组成$Q_u$，得到效用保护投影矩阵$P_u=Q_uQ_u^{\mathsf T}$。$\rho_u$或$r_u$控制需要显式保护的效用梯度覆盖范围。

若一次参数更新沿着$Q_u$所张成的方向移动，可能改变无害任务损失的一阶近似值。因此，本文将$\|Q_u^{\mathsf T}\delta\|_2^2$作为对更新在效用敏感方向分量的惩罚。安全子空间与效用保护子空间可能重叠，本文不假设二者天然正交；若强制将安全方向完全投影到效用子空间的正交补，可能同时消除部分必要的安全更新。本文通过可调惩罚项和无害样本上的输出保持损失控制这种冲突，并在实验中报告安全与效用的共同变化。

梯度子空间是对有限样本和初始参数邻域的近似。为避免将训练样本兼作最终效用证据，$\mathcal{D}_{u}^{\mathrm{train}}$仅用于基构造和优化，$\mathcal{T}_{u}$用于独立评估。无害样本需覆盖与编辑目标相近但本身合理的请求，以观察方法是否把语义相近误当作风险信号。

\subsection{双子空间约束机制}

本文将编辑更新参数化为安全子空间基上的线性组合：
\begin{equation}
  \delta=Q_s z,\qquad z\in\mathbb{R}^{r_s}.
  \label{eq:update-param}
\end{equation}
式\eqref{eq:update-param}限制编辑更新位于安全梯度所张成的方向范围内，避免优化器在无约束参数空间中产生与安全训练目标无关的任意更新。对效用保护子空间的约束采用软惩罚：
\begin{equation}
  \mathcal{R}_{\mathrm{sub}}(\delta)
  =\|Q_u^{\mathsf T}\delta\|_2^2
  =\|Q_u^{\mathsf T}Q_s z\|_2^2.
  \label{eq:utility-subspace-penalty}
\end{equation}
相较于硬正交约束，软惩罚能够处理$Q_s$和$Q_u$存在交叠的情况；其权重$\lambda_{\mathrm{sub}}$决定效用保护方向与安全编辑目标之间的折中。

为从行为层面补充参数子空间约束，定义无害任务上的教师分布保持项：
\begin{equation}
  \mathcal{R}_{\mathrm{KL}}(\delta)
  =\frac{1}{|\mathcal{D}_{u}^{\mathrm{train}}|}
  \sum_{x\in\mathcal{D}_{u}^{\mathrm{train}}}
  D_{\mathrm{KL}}\!\left(
  p_{\theta_0}(\cdot\mid x)
  \,\Vert\,
  p_{\theta_0+\delta}(\cdot\mid x)
  \right).
  \label{eq:utility-kd}
\end{equation}
该项以编辑前模型在无害输入上的输出分布作为参照，减少不必要的行为漂移。若任务具有标准答案，也可同时保留任务损失；具体组合应在实验配置中统一说明。仅有KL保持项并不代表任务效用必然提高，因此最终效用仍由独立测试集衡量。

不考虑后续微调时的双子空间目标写为
\begin{equation}
  \min_{z}\quad
  \mathcal{L}_{s}(\theta_0+Q_s z;\mathcal{D}_{s}^{\mathrm{train}})
  +\lambda_{\mathrm{sub}}\mathcal{R}_{\mathrm{sub}}(Q_s z)
  +\lambda_{\mathrm{KL}}\mathcal{R}_{\mathrm{KL}}(Q_s z)
  +\lambda_{2}\|Q_s z\|_2^2 .
  \label{eq:dual-subspace-objective}
\end{equation}
其中$\lambda_{\mathrm{KL}}$控制输出保持强度，$\lambda_2$限制整体更新幅度。式\eqref{eq:dual-subspace-objective}提供当前编辑效果与无害效用保持的共同优化基础，下一小节再将其扩展到代表性微调扰动下的鲁棒目标。

\subsection{微调扰动建模与鲁棒优化目标}

在编辑前，本文利用训练用微调环境集合$\mathcal{E}_{\mathrm{train}}$获得扰动模板库$\mathcal{B}_{\mathrm{train}}=\{\Delta_{0}^{(0)},\Delta_{0}^{(1)},\ldots,\Delta_{0}^{(M)}\}$。每个非零模板由基础模型在一组指定下游数据和固定训练预算下微调得到。编辑目标中的安全损失在$\theta_0+\delta+\Delta_{0}^{(j)}$处计算，作为对后续任务更新方向的近似训练信号。该加法模板不等同于在编辑模型上重新微调，因此完整实验还需从编辑模型实际执行微调，直接测量微调后的安全表现。训练扰动环境与最终评测环境应按任务或数据源划分；若两者完全相同，结果只能说明对已见微调环境的适应，不能据此声称泛化到未见微调任务。

为同时考虑扰动库内的平均表现和较差环境，定义平滑最坏情形安全损失：
\begin{equation}
  \mathcal{L}_{\mathrm{rob}}(z)
  =(1-\alpha)\frac{1}{M+1}
    \sum_{j=0}^{M}\mathcal{L}_{s}
    \!\left(\theta_0+Q_s z+\Delta_{0}^{(j)};\mathcal{D}_{s}^{\mathrm{train}}\right)
  +\alpha\tau\log\!\left[
    \frac{1}{M+1}
    \sum_{j=0}^{M}
    \exp\!\left(
      \frac{\mathcal{L}_{s}(\theta_0+Q_s z+\Delta_{0}^{(j)};
      \mathcal{D}_{s}^{\mathrm{train}})}{\tau}
    \right)
  \right],
  \label{eq:robust-loss}
\end{equation}
其中$\alpha\in[0,1]$控制最坏环境项的权重，$\tau>0$为平滑温度。当$\alpha$增大或$\tau$减小时，目标更强调扰动库中安全损失较大的微调环境；$\Delta_{0}^{(0)}=\mathbf{0}$使目标同时覆盖不微调时的编辑效果。该目标对有限扰动库进行优化，不构成对任意参数扰动的极小极大保证。

最终优化问题为
\begin{equation}
  \min_{z}\quad
  \mathcal{J}(z)
  =\mathcal{L}_{\mathrm{rob}}(z)
  +\lambda_{\mathrm{sub}}\|Q_u^{\mathsf T}Q_s z\|_2^2
  +\lambda_{\mathrm{KL}}\mathcal{R}_{\mathrm{KL}}(Q_s z)
  +\lambda_{2}\|Q_s z\|_2^2 .
  \label{eq:full-objective}
\end{equation}
式\eqref{eq:full-objective}包含三类互补约束：安全子空间限定可行编辑方向，效用保护子空间与KL项限制对无害任务的干扰，扰动库上的鲁棒损失降低编辑对已建模后续微调环境的敏感性。各项权重通过验证集确定，并在实验中报告其敏感性。

需要强调的是，本方法不假设后续微调扰动可被完全预测。扰动库仅提供可重复的训练环境近似；对未见环境的鲁棒性必须通过留出的微调任务独立检验。与CKU等训练式安全知识遗忘方法、ILU等面向下游微调鲁棒性的训练方法，以及基于Fisher信息与功能干预约束的FDCU相比，本文研究参数级安全编辑，并将编辑更新限制在安全梯度子空间内；鲁棒项使用基础模型在训练用微调环境上获得的参数差作为虚拟扰动模板，在模板叠加下计算安全损失，最终通过从编辑模型实际微调的独立测试检验该近似。与针对保留知识进行零空间投影的编辑方法相比，本文同时以安全目标梯度和无害任务梯度构造两个子空间，并增加后续微调环境维度的评价\cite{fang2024alphaedit,wang2025ilu,shi2025cku,li2026fdcu}。该差异是本文的设计定位，是否构成充分的新颖性仍需结合完整相关工作和实验对比论证，不在本稿中使用“首次”等绝对表述。

\subsection{算法流程与复杂度分析}

双子空间鲁棒编辑流程概述如下。实际实现时，编辑参数集合$\Theta_e$应在实验前固定；梯度基、扰动库和编辑参数必须只由训练数据构建，验证集用于选择维度和损失权重，测试集仅用于最终报告。

\begin{enumerate}
  \item \textbf{输入：}初始模型$\theta_0$、安全训练集$\mathcal{D}_{s}^{\mathrm{train}}$、无害训练集$\mathcal{D}_{u}^{\mathrm{train}}$、微调训练环境集合$\mathcal{E}_{\mathrm{train}}$及超参数。
  \item \textbf{构造安全子空间：}计算安全样本对$\Theta_e$的梯度矩阵$G_s$，按能量阈值或验证集选择秩$r_s$，得到基$Q_s$。
  \item \textbf{构造效用保护子空间：}计算无害训练样本梯度矩阵$G_u$，得到基$Q_u$及对应秩$r_u$。
  \item \textbf{构造扰动库：}在各微调训练环境上微调$\theta_0$，记录参数差模板$\Delta_{0}^{(j)}$，并加入零扰动$\Delta_{0}^{(0)}$。
  \item \textbf{优化编辑系数：}初始化$z=0$，根据式\eqref{eq:full-objective}计算损失和梯度，使用选定优化器迭代更新$z$；对扰动环境进行小批量采样或分批梯度累积。
  \item \textbf{输出：}令$\delta^*=Q_s z^*$，得到编辑模型$\theta_e=\theta_0+\delta^*$；在独立安全、无害效用和微调测试集上进行评估。
\end{enumerate}

若编辑参数维度为$d_e$，安全和无害梯度样本数分别为$N_s$与$N_u$，保留子空间维度为$r_s$与$r_u$，则显式保存两个稠密基的空间开销为$O(d_e(r_s+r_u))$，扰动库的存储开销与保存的微调参数差数量及其压缩表示有关。梯度矩阵的精确SVD可能成为主要计算开销；实现时可采用分块梯度累积、随机截断SVD或在各选定层内分别构造低秩基，以减少峰值显存。每次编辑优化需对抽取的微调环境执行安全损失前向和反向计算；若每步使用$b$个扰动环境，额外前向和反向计算量随$b$近似线性增加。复杂度的具体数值需在实现模型、编辑层和硬件确定后实测。

\section{实验结果与分析}

本节给出用于验证方法的实验设计和待填结果位置。当前尚未执行相关实验，以下内容不代表实测结果。正式写作时应以实际运行配置、日志和结果表替换待补项。

\subsection{实验设置}

\textbf{模型与编辑参数。}项目现有实验计划中列有Llama-2-7B-Chat、Qwen-7B和GPT-J-6B等候选模型。正式实验应选定至少一个主模型，并根据编辑器兼容性与计算资源决定是否增加跨模型验证；需记录模型版本、分词器、聊天模板、精度格式及可训练参数范围。若不同编辑器支持的层或参数不同，应统一评价协议并说明差异。

\textbf{安全与无害数据。}安全编辑训练集和测试集可围绕SafeEdit所覆盖的不安全类别构建，并按原始请求、攻击改写和语义相近提示进行分层划分\cite{wang2024safeedit}。无害训练样本用于构建$Q_u$及KL保持项，无害测试样本用于独立评估一般任务能力和与风险语义相近的正常请求。具体数据来源、样本规模、类别比例、去重规则、标签说明和训练/验证/测试划分均需在实验前固定并记录。安全测试提示不得参与梯度基构造或超参数拟合。

\textbf{对比方法。}建议至少比较未编辑模型、直接安全编辑基线及在相同目标数据上适配的参数级编辑器。现有计划列出的ROME、MEMIT和AlphaEdit可作为候选编辑基线；CKU等训练式安全遗忘方法可作为补充参照，但需将编辑范式和训练成本单独说明，避免把不同算法类别混为同一基线\cite{meng2022rome,meng2023memit,fang2024alphaedit,shi2025cku}。本方法还需设置仅双子空间约束、仅扰动鲁棒优化和完整方法等变体。

\textbf{后续微调环境。}选择若干与安全编辑数据不重叠的下游任务，预先固定微调方法与训练预算。扰动库使用训练环境生成；至少保留一个数据源或任务类型作为未见环境测试。报告全参数微调或LoRA等方式、数据量、训练步数、学习率、批大小和随机种子。所有编辑方法都使用相同的微调与测试配置。

\textbf{报告规范。}对随机过程采用多个随机种子重复运行，报告均值与标准差；自动安全判别器须说明版本，并对一定比例的输出进行人工复核。报告基础模型、编辑模型和各类微调后模型的原始分数，不仅报告相对提升。若不同模型无法使用完全相同的编辑器，应分别提供同模型内比较和跨模型趋势。

表\ref{tab:ch2-config-todo}列出当前需在正式实验前确定的配置。
\begin{table}[htbp]
  \centering
  \caption{第二章实验配置待确定项}
  \label{tab:ch2-config-todo}
  \begin{tabular}{p{0.25\textwidth}p{0.58\textwidth}}
    \hline
    配置项 & 待确定内容 \\ \hline
    主模型与编辑层 & 模型版本、可用编辑器、$\Theta_e$范围 \\
    安全数据 & 训练/验证/测试样本、风险类别、攻击改写划分 \\
    无害数据 & 子空间构造数据、KL保持数据、独立效用测试集 \\
    微调环境 & 训练环境、留出环境、微调算法与训练预算 \\
    超参数 & $r_s,r_u,\rho_s,\rho_u,\alpha,\tau,\lambda_{\mathrm{sub}},\lambda_{\mathrm{KL}},\lambda_2$ \\
    运行设置 & 硬件、精度、随机种子、生成和评估配置 \\
    \hline
  \end{tabular}
\end{table}

\subsection{安全编辑效果对比}

本实验在尚未进行后续微调的条件下，比较各方法的即时安全编辑效果与无害效用。对每种方法报告安全测试集上的ASR、SRR，以及无害任务准确率或评分、无害相邻请求误拒率和编辑局部性。对于风险类型和提示来源，还应给出分类别结果，检验总体均值是否掩盖特定类别的失败。

主要结果表建议包含方法、编辑后ASR、SRR、无害效用、误拒率及运行时间等列；必要时另列编辑参数量和训练/编辑步数。现阶段数据未运行，相关数值留待实验后填写。分析应指出安全效果与效用之间的实际变化，并结合样例检查是否存在以空回答、重复拒答或明显降质换取较低ASR的伪安全现象。

% TODO: 插入各编辑器的即时安全效果与无害效用对比表；补充按风险类别的结果和输出案例。

\subsection{不同微调策略下的安全编辑鲁棒性分析}

本实验在编辑完成后，从编辑模型本身重新执行扰动库内微调及留出的微调任务，再使用相同安全测试集评价编辑保持情况。具体地，对每种微调环境计算$\widehat{\theta}_{\mathrm{ft}}^{(j)}=\mathrm{FT}_{e_j}(\theta_0+\delta)$，并在该模型上测量安全指标。该步骤评估真实的编辑后微调过程，与训练阶段使用固定扰动模板的虚拟目标相区分。至少报告编辑后、微调后绝对ASR/SRR，保持率$\mathrm{Ret}^{(j)}$，以及对应下游任务效用。可按微调方式、任务类型和训练预算分组，观察微调方法与编辑方法之间是否存在交互。

为验证扰动建模的泛化能力，应将训练扰动环境与最终测试环境分开。若完整方法只在训练扰动库对应的微调任务上表现较好，而在留出任务上没有相同趋势，应据实限定结论。可通过不同训练步数的曲线展示安全指标变化，但曲线采样点应在实验前设定，并保持各方法一致。

% TODO: 插入不同微调环境的ASR/SRR与保持率结果；标明训练环境和留出环境。

\subsection{双子空间约束消融实验}

本实验用于分别验证安全子空间和效用保护子空间的作用。建议比较以下变体：完整双子空间方法；移除安全子空间参数化、在相同编辑参数范围内直接优化；移除效用保护子空间惩罚；同时移除效用子空间惩罚和KL保持项。其余数据、微调环境和优化预算保持一致。评价指标至少包括即时SRR、无害效用、无害相邻请求误拒率及微调后的SRR。

若移除安全子空间后更新幅度增大但安全目标无相应收益，或无害效用明显下降，可分析其是否与更新方向缺少约束有关；若移除效用保护后安全指标改善而一般效用下降，则应报告该权衡，而不将单一安全指标作为整体优劣依据。以上均为待验证假设，不应在得到实验数据前写成结果。

% TODO: 插入双子空间相关消融表，区分安全子空间、效用子空间和行为保持项。

\subsection{鲁棒优化消融实验}

本实验用于检验扰动库与平滑最坏情形项的作用。建议比较不含微调扰动的双子空间编辑、仅在扰动环境上优化平均安全损失、以及采用式\eqref{eq:robust-loss}的完整目标。为避免计算预算成为混杂因素，应尽量使各变体的编辑迭代次数、训练数据访问量和参数选择流程一致。

除报告扰动库内的保持率外，还需单独报告留出微调任务的结果。若完整目标提升已见环境表现但未提升留出环境表现，应讨论扰动库覆盖不足或环境差异，而不能笼统宣称具有微调鲁棒性。必要时可比较不同扰动库规模和环境多样性。

% TODO: 插入平均目标、平滑最坏情形目标和无扰动目标的消融结果。

\subsection{通用能力与知识局部性分析}

本实验评估安全编辑是否损害与目标风险无关的模型行为。通用能力测试应根据模型的指令跟随、问答或其他主要用途选取任务，并报告编辑前后的任务级结果。知识局部性测试应包括与编辑目标无关的问题、语义相近但无害的问题，以及必要的同义改写请求；分别报告输出分布差异、任务性能变化和误拒率。

若方法仅在整体效用均值上保持稳定，但对相近无害请求的误拒明显上升，应在结果中单独指出。分析还应比较效用子空间惩罚和KL项对不同类别无害任务的影响，说明保护方向是否覆盖了评测数据中的主要变化。

% TODO: 插入通用任务、无关知识和相邻无害请求的分项结果。

\subsection{参数敏感性分析}

参数敏感性实验重点考察安全子空间秩$r_s$、效用保护子空间秩$r_u$、能量保留阈值$\rho_s$与$\rho_u$、扰动环境权重$\alpha$、平滑温度$\tau$及损失权重$\lambda_{\mathrm{sub}}$、$\lambda_{\mathrm{KL}}$和$\lambda_2$。每次改变一个参数时固定其余设置；对计算成本较高的参数，可先在验证集缩小范围，再在测试集只报告选定配置。

参数分析应同时观察安全效果、微调后保持、无害任务效用和编辑计算开销。若某个参数只改善一种指标，却导致另一关键指标明显下降，应呈现权衡曲线或多指标表格，不应只报告最优单点。最终参数应依据验证集确定，并在测试集上一次性评估。

% TODO: 插入主要参数的敏感性曲线或表格，并记录参数选择规则。

\section{本章小结}

本章围绕安全编辑在后续微调中的保持问题，给出双子空间约束与扰动鲁棒优化的方案设计。方法首先基于安全目标梯度构造安全子空间，再依据无害任务梯度构造效用保护子空间；通过安全子空间参数化、效用方向惩罚和无害输出保持项限制编辑更新，并在代表性后续微调扰动库上优化安全目标。该方法的适用范围由编辑参数集合、梯度样本覆盖和微调环境库共同决定，不能据此推出对任意微调或任意风险输入的保证。

本章还给出了安全编辑效果、微调后保持、无害效用、局部性和参数敏感性等实验设计。由于实验尚未完成，本章不对方法是否有效或相较基线提升多少作结论。待完成实验后，应根据实测结果补充预实验分析、各项对比与消融表格，并据此修订摘要、绪论和本章小结中的效果性表述。
